\section{行业背景分析}
\subsection{行业现状}
智慧家电、工业控制、农业监测与安防等场景，都需要把现场数据转化为判断或控制依据。将模型推理部署在嵌入式端，可以使数据在设备附近完成处理，再按需要输出结果或与外部系统协作。企业命题将感知、通信和智能计算放入同一系统，关注模型与设备、操作系统及数据接口的结合。\cite{challenge}

嵌入式推理开发面临两类相互关联的问题。一类是计算映射：矩阵运算与偏置、激活、数据整理等步骤具有不同结构，需要选择合适的处理单元。另一类是工程衔接：模型的数据类型、量化规则、缓冲区和调用方式，需要在转换器、运行时与系统软件之间保持一致。软件体系决定加速器的计算能力能否用于完整模型。\cite{gemmini,technicalreport}

开放硬件使处理器与加速器内部结构可见，为算子优化和接口验证提供条件；可配置设计也意味着硬件参数变化会影响软件接口。因此，模型部署需要同时检查数值正确性、硬件配置和系统适配。本项目以模型与算子为中心，将计算后端、部署方式和验证记录放在统一接口下管理。

\subsection{命题相关技术调研}
在通用与向量计算方面，\ProcessorCore{}承担操作系统、控制程序和调度，\VectorUnit{}与其连接，提供 RVV 1.0 向量执行能力。向量指令适合表达偏置、激活、重定标和数据整理等多元素操作，能够与矩阵加速形成分工。软件需要使用匹配的编译目标与执行环境，确保算子实际调用向量指令。\cite{rocket,saturn,rvv,gcc}

在矩阵计算方面，\AcceleratorName{}提供可配置计算阵列及数据搬运接口，通过 RoCC 自定义指令接收处理器任务。其 C 接口可作为矩阵后端的基础；模型转换、算子注册、后端分配和数值验证则由项目软件层衔接。\cite{gemmini}

在硬件集成方面，\IntegrationFramework{}组织处理器、加速器、互连和外设，\MemoryController{}提供开放的 DRAM 控制器与 PHY 生成基础。已有组件和接口为集成提供依据，但框架并非本试验箱的现成支持包，板级约束、存储训练及组合后的资源与时序仍需验证。\cite{chipyard,litedram,technicalreport}

在操作系统方面，\RealTimeSystem{}与\LinuxDistribution{}分别提供实时系统和可裁剪应用环境。两者共用模型、算子核心源码和硬件接口，分别构建镜像与程序，通过切换运行环境开展适配。这样可以把工作集中于板级支持和计算后端衔接，不额外引入双系统并发。\cite{rtthread,linux}

本方案据此采用开源组件作为计算与系统基础，把新增研究集中于统一算子描述、受限模型转换、三后端映射和逐层验证。选型依据与固定源码记录纳入\TechnicalBaseline{}技术报告，后续以同一模型、同一数值约定和可重复测试评估系统可用性。\cite{technicalreport}
